\section{Background}\label{sec:background}

This section recalls the cryptographic primitives used by \thunderbolt. All notation introduced here is used without further comment in subsequent sections.

\subsection{Elliptic-Curve Groups}\label{sec:bg-ec}

Let $\GG$ be an elliptic-curve group of prime order $q$ with generator $G$ (or $H$, when a second independent base point is needed). Bitcoin uses the secp256k1 curve; \thunderbolt works over the same group and order. A \emph{discrete-log pair} is a pair $(x, X)$ with $X = xH$ for some scalar $x \in \ZZ_q$. Given only $X$ and $H$, no efficient algorithm can recover $x$---this is the \emph{elliptic-curve discrete-logarithm problem} (ECDLP), on which the security of Bitcoin and of our proofs rests.

\subsection{The Paillier Cryptosystem}\label{sec:bg-paillier}

The Paillier cryptosystem~\cite{paillier1999public} is an asymmetric encryption scheme whose plaintext space is $\ZZ_N$ for an RSA modulus $N = pq'$ (we write $q'$ to distinguish the Paillier factor from the curve order $q$). The public key is $(N, g)$ with $g \in \ZZ_{N^2}^*$ of order divisible by $N$. Encryption of a message $m \in \ZZ_N$ with randomness $\rho \in \ZZ_N^*$ is
$$
\Enc_{pk}(m;\rho) = g^m \rho^N \bmod N^2.
$$
Decryption uses the factorization of $N$. The scheme is semantically secure under the \emph{Decisional Composite Residuosity} (DCR) assumption.

\paragraph{Additive homomorphism.} Paillier is \emph{additively} homomorphic: given ciphertexts $c_1 = \Enc(m_1)$ and $c_2 = \Enc(m_2)$, the product
$$
c_1 \cdot c_2 \bmod N^2 = \Enc(m_1 + m_2)
$$
is a valid encryption of $m_1 + m_2 \bmod N$. We write this operation as $c_1 \oplus c_2$, and the corresponding ``ciphertext subtraction'' as $c_1 \ominus c_2 := c_1 \cdot c_2^{-1} \bmod N^2$. This additive homomorphism is the reason \thunderbolt can compose a fresh holder ciphertext without decrypting the invoice ciphertext.

\subsection{Sigma Protocols and Fiat--Shamir}\label{sec:bg-sigma}

A \emph{Sigma protocol} is a three-move honest-verifier zero-knowledge proof for a relation $\mathcal{R}$: the prover sends a commitment $A$, receives a challenge $c$ from the verifier, and responds with $z$. The verifier accepts iff a public verification equation involving $(A, c, z)$ and the statement holds. Sigma protocols are complete, special-sound, and honest-verifier zero-knowledge.

The \emph{Fiat--Shamir heuristic}~\cite{fiat1987prove} turns a Sigma protocol into a non-interactive proof (NIZK) in the random-oracle model by replacing the verifier's challenge with $c = \Hash(\text{statement}, A)$, where $\Hash$ is a cryptographic hash function modeled as a random oracle. The resulting proof $\pi = (A, z)$ (or, equivalently, $(c, z)$ with the commitment recomputed) is publicly verifiable.

\subsection{Verifiable Encryption of Discrete Logarithms}\label{sec:bg-vedl}

Camenisch and Shoup~\cite{camenisch2003verifiable} give an efficient protocol for \emph{verifiable encryption of discrete logarithms}: given a discrete-log pair $(m, M = mH)$ and a Paillier ciphertext $c = \Enc_{pk}(m)$, the prover produces a Sigma-protocol proof that $c$ encrypts the same integer whose discrete log is $M$. The proof relation is
$$
\mathcal{R}_{\mathrm{VE}} = \{(M, c;\; m, \rho) : M = mH \;\land\; c = g^m \rho^N \bmod N^2\}.
$$

The key technique is a \emph{shared response}: the Sigma-protocol response $z$ simultaneously satisfies both the elliptic-curve verification equation ($zH \stackrel{?}{=} A_{\mathrm{EC}} + cM$) and the Paillier verification equation ($g^z\tau^N \stackrel{?}{=} A_{\mathrm{Pai}} \cdot c^c$), where the same $z$ forces the same underlying integer in both relations. This ``shared-$z$'' construction is the technical core of our proofs.

\subsection{Verifiable Encrypted Secret Sharing (VESS)}\label{sec:bg-vess}

Shoup's \emph{Simple VESS}~\cite{shoup2024simple} provides a mechanism for verifiably encrypting a Shamir secret sharing to a committee of $n$ shareholders. The sharing party (the ``dealer'') publishes encrypted shares $(\hat{s}_1, \ldots, \hat{s}_n)$ together with a proof that there exists a unique polynomial of degree at most $t-1$ whose evaluations are the plaintexts of these ciphertexts. Any observer can verify the consistency of the encrypted shares without decrypting them. If at least $t$ of the $n$ shareholders are honest, they can collectively reconstruct the shared secret.

In \thunderbolt, VESS is used when the Committee is instantiated as a threshold committee rather than a single actor. The committee-side refresh value $r_s$ is shared via a VESS protocol so that the compensating update $y_{s+1} = y_s - r_s$ is performed collectively, without any single committee member learning $y_s$ or $r_s$ in the clear. This modular substitution replaces the single-Committee trust assumption with a standard threshold-honesty assumption, and it can be performed without any change to Core Proofs~1 and~2.
